\section{Introduction}
\label{sec:introduction}

The source panel method is a numerical technique used in aerodynamics to approximate the potential flow around bodies of
arbitrary shape.
It belongs to the family of boundary element methods, where the surface of a body is discretized into a finite number of
panels.
Each panel is assigned a source distribution whose strength is determined such that the flow satisfies the
impermeability boundary condition onSource Panel Method the body surface.
The method is based on the assumptions of incompressible, irrotational and inviscid flow, allowing the velocity field to
be represented by a scalar potential that satisfies Laplace's equation.
By superimposing the effects of the source panels and the free-stream flow, the velocity potential around the body can
be computed.
Once the velocity distribution is known, aerodynamic quantities such as the pressure coefficient can be obtained using
Bernoulli's equation.
The source panel method is particularly useful for predicting pressure distributions on non-lifting bodies and serves as
a foundation for more advanced panel methods that incorporate vortices to model lifting flows.